\section{Minimize the Maximum Distance Between Gas Stations}
\label{sec:bs-gas-stations}

\problemstatement{We are given $n$ gas station positions, sorted along a
number line, and an integer $k$. We may insert $k$ additional gas stations
\textbf{anywhere} (not restricted to integer positions) among the existing
ones. We must choose where to insert them so as to \textbf{minimize the
largest gap} between any two consecutive gas stations (existing or newly
inserted), and report that minimized largest gap.}

\begin{patternbox}
This problem keeps the ``minimize the maximum, subject to a budget''
shape of Sections~\ref{sec:bs-book-allocation}--\ref{sec:bs-painters-partition},
but with one crucial new wrinkle: the answer is a \textbf{real number},
not an integer, because stations may be placed at fractional positions.
Whenever a binary-search-on-the-answer problem allows continuous
(real-valued) placement rather than discrete/integer choices, the
\texttt{while(lo <= hi)} integer loop must be replaced with a
fixed-precision loop that runs until \textit{hi} and \textit{lo} are
within some small tolerance $\varepsilon$ of each other.
\end{patternbox}

\subsection*{Understanding the problem carefully}

Consider the existing gaps between consecutive original stations:
$\textit{gap}_1, \textit{gap}_2, \dots, \textit{gap}_{n-1}$. For a fixed
candidate maximum allowed gap $d$, the minimum number of new stations
needed to bring gap $\textit{gap}_i$ down to segments no larger than $d$
is
\[
  \left\lceil \frac{\textit{gap}_i}{d} \right\rceil - 1
\]
(inserting $c$ points into a single gap of length $L$ divides it into
$c+1$ equal-or-smaller segments; we need the smallest $c$ such that
$L/(c+1) \le d$, which rearranges to $c \ge L/d - 1$, i.e.\ $c = \lceil
L/d \rceil - 1$ for integer $c$). Summing this over every existing gap
gives $\textit{stationsNeeded}(d)$, the total new stations required to
guarantee no gap exceeds $d$. As $d$ increases, fewer stations are needed
per gap, so $\textit{stationsNeeded}(d)$ is non-increasing in $d$.

\begin{ideabox}[Candidate Space and Predicate]
The candidate space is real numbers $d \in (0, \max(\textit{gap}_i)]$.
The predicate is $P(d) = (\textit{stationsNeeded}(d) \le k)$, and we want
the smallest $d$ where $P$ holds.
\end{ideabox}

\subsection*{Why a fixed-precision loop replaces the integer loop}

With a continuous candidate space, there is no notion of ``the next
integer candidate,'' so the loop cannot terminate by \textit{lo} and
\textit{hi} crossing as integers do. Instead, we run the loop a fixed
number of times (or until $\textit{hi} - \textit{lo}$ shrinks below a
chosen precision threshold $\varepsilon$, e.g.\ $10^{-6}$), narrowing
\textit{lo} or \textit{hi} to \textit{mid} (not $\textit{mid} \pm 1$, since
there is no discrete step to add or subtract) at each iteration. Each
iteration still halves the remaining interval, so a fixed number of
iterations (around $50$--$100$, since $2^{-100}$ is far smaller than any
reasonable $\varepsilon$) is always enough regardless of the initial
interval width.

\subsection*{Algorithm}

\begin{enumerate}
  \item Compute the existing gaps; $\textit{lo} \gets 0$,
        $\textit{hi} \gets \max(\textit{gap}_i)$.
  \item Repeat a fixed number of times (or while
        $\textit{hi} - \textit{lo} > \varepsilon$):
  \begin{enumerate}
    \item $\textit{mid} \gets (\textit{lo}+\textit{hi})/2$.
    \item Compute $\textit{stationsNeeded}(\textit{mid}) =
          \sum_i \left( \lceil \textit{gap}_i / \textit{mid} \rceil - 1
          \right)$.
    \item If $\textit{stationsNeeded}(\textit{mid}) \le k$:
          $\textit{hi} \gets \textit{mid}$ (feasible; try smaller).
    \item Else: $\textit{lo} \gets \textit{mid}$ (infeasible; try larger).
  \end{enumerate}
  \item Return $\textit{hi}$ (or $\textit{lo}$; they have converged to
        within $\varepsilon$).
\end{enumerate}

\begin{algorithm}[H]
\caption{Minimize Maximum Gas Station Distance}
\begin{algorithmic}[1]
\Procedure{MinimizeMaxDistance}{\textit{stations}, $k$}
  \State compute gaps $\textit{gap}_1, \dots, \textit{gap}_{n-1}$
  \State $\textit{lo} \gets 0$, $\textit{hi} \gets \max_i \textit{gap}_i$
  \For{$\textit{iter} \gets 1$ to $100$}
    \State $\textit{mid} \gets (\textit{lo}+\textit{hi})/2$
    \State $\textit{needed} \gets \sum_i (\lceil \textit{gap}_i / \textit{mid} \rceil - 1)$
    \If{$\textit{needed} \le k$}
      \State $\textit{hi} \gets \textit{mid}$
    \Else
      \State $\textit{lo} \gets \textit{mid}$
    \EndIf
  \EndFor
  \State \Return \textit{hi}
\EndProcedure
\end{algorithmic}
\end{algorithm}

\subsection*{Dry run}

\dryrun{Take stations at $[1, 2, 3, 4, 5]$ (gaps all $=1$), $k = 4$.}

There are $4$ gaps, each of length $1$. With $k=4$ new stations to
distribute (one per gap), each gap of length $1$ becomes two segments of
length $0.5$. So $\textit{stationsNeeded}(0.5) = \sum (\lceil 1/0.5
\rceil - 1) = \sum (2-1) = 4 \cdot 1 = 4 \le 4$ -- feasible. Checking a
slightly smaller candidate, say $d=0.49$:
$\lceil 1/0.49 \rceil - 1 = \lceil 2.04 \rceil - 1 = 3-1=2$ per gap, times
$4$ gaps $=8 > 4$ -- infeasible. So the binary search converges to
$d = 0.5$, the minimized maximum distance.

\subsection*{C++ implementation}

\begin{lstlisting}
#include <bits/stdc++.h>
using namespace std;

double minimizeMaxDistance(vector<int>& stations, int k) {
    int n = stations.size();
    vector<double> gaps(n - 1);
    for (int i = 0; i + 1 < n; i++) {
        gaps[i] = stations[i + 1] - stations[i];
    }

    double lo = 0, hi = *max_element(gaps.begin(), gaps.end());

    for (int iter = 0; iter < 100; iter++) {
        double mid = (lo + hi) / 2.0;
        long long needed = 0;
        for (double g : gaps) {
            needed += (long long)ceil(g / mid) - 1;
        }

        if (needed <= k) {
            hi = mid;
        } else {
            lo = mid;
        }
    }
    return hi;
}

int main() {
    vector<int> stations = {1, 2, 3, 4, 5};
    cout << fixed << setprecision(4);
    cout << "Minimized maximum distance: "
         << minimizeMaxDistance(stations, 4) << endl;
    return 0;
}
\end{lstlisting}

\begin{complexitybox}
\textbf{Time Complexity.} The fixed-precision loop runs a constant number
of iterations (e.g.\ $100$), and each feasibility check costs $O(n)$,
giving
\[
  O(n \cdot \text{iterations}) = O(n),
\]
treating the iteration count as a constant determined by the desired
precision (more precisely, $O(n \log(\textit{range}/\varepsilon))$ if the
precision $\varepsilon$ is treated as a variable rather than fixed).
\textbf{Space Complexity.} $O(n)$ to store the gaps.
\end{complexitybox}

\begin{takeawaybox}
Binary search on the answer is not limited to integer candidates. When the
problem allows real-valued placement, keep every other part of the
template (monotonic feasibility predicate, minimize/maximize narrowing
direction) and simply replace the integer loop and its $\pm 1$ adjustments
with a fixed-iteration or epsilon-tolerance loop that narrows
\textit{lo}/\textit{hi} directly to \textit{mid}.
\end{takeawaybox}
